% gt_part_L.tex -- [GT] Part L: the literature register (L1).  Slot 16, package L16-3.
% Read by lean_census.py (gt_part_*.tex).  Each entry is a hypothesis structure of GL2/Literature/*.lean:
% a Prop-valued structure, never an axiom.  Quotations are verbatim from the library scans (read-only);
% "UNVERIFIED / NOT IN LIBRARY" marks a locator that has not been read.
\newtheorem{hypothesis}[theorem]{Literature hypothesis}

\chapter{Part L: the literature register}

\section{Modular forms modulo a prime}

\begin{hypothesis}[Serre, Bourbaki 416, \S1.2, Thm.~2]\label{lit:serre416-thm2}
  \lean{GL2.Literature.Bourbaki416Thm2} \leanok
  \textbf{Locator: verified (English translation, p.~4; the French original in the library is a truncated PDF).}
  ``Theorem 2. Suppose $f$ and $f'$ are $p$-integral modular forms of weights $k$ and $k'$. If $f\equiv f'\pmod p$,
  and $f\not\equiv0\pmod p$, then one has $k\equiv k'\pmod{p-1}$.'' Standing hypothesis $p\ge5$.
  Same statement: Swinnerton-Dyer, LNM 350, p.~22, Thm.~2(iv).
\end{hypothesis}

\begin{hypothesis}[Serre, Bourbaki 416, Remark after Thm.~2]\label{lit:serre416-pow}
  \lean{GL2.Literature.Bourbaki416RemarkPow} \leanok
  \textbf{Locator: verified as stated (translation p.~4), stated without proof by Serre.}
  ``Under the hypotheses of Theorem 2, if $f\equiv f'\pmod{p^n}$, then one can show that
  $k\equiv k'\pmod{p^{n-1}(p-1)}$.'' [E] cites this as Serre LNM 350, \S1, Th\'eor\`eme~1: \textbf{UNVERIFIED,
  NOT IN LIBRARY} (LNM 350 pp.~191--268 is absent).
\end{hypothesis}

\begin{hypothesis}[Sturm's bound, level one]\label{lit:sturm}
  \lean{GL2.Literature.SturmLevelOne} \leanok
  \textbf{Locator: verified (statement and page), literature-list pass 2026-10-05, LIT-13}: Sturm, \emph{On the congruence of
  modular forms}, LNM 1240 (1987), pp.~275--280, Thm.~1, p.~276 (printed $=$ PDF $+\,274$): $\Gamma\supseteq\Gamma(N)$, $k\ge1$,
  $f,g\in M_k(\Gamma,R_F)$, $\operatorname{ord}_\lambda(f-g)>k[\Gamma(1):\Gamma]/12\Rightarrow f\equiv g\bmod\lambda$; at level one the
  bound is $k/12$, the constant term included, as in \texttt{SturmLevelOne}. Sturm's coefficient ring is $R_F$ and the prime
  $\lambda$; the passage to $p$-integral forms is a localisation not written there. The proof uses Shimura Thm.~6.6 (Ch.~6, not
  held), but the level-one case is by $\Delta$ and $j$. No quotation.
\end{hypothesis}

\begin{hypothesis}[Bellaïche--Khare, Thms.~I--III]\label{lit:bk}
  \lean{GL2.Literature.BellaicheKhareThms} \leanok
  \textbf{Locator: verified, Compositio Math.\ 151 (2015), p.~401.}
  ``Theorem I. Assume that $\bar\rho$ is unobstructed. Then the morphism $\tilde R^0_{\bar\rho}\to A_{\bar\rho}$ is an
  isomorphism, and $A_{\bar\rho}$ is isomorphic to a power series ring in two variables $F[[x,y]]$.'' Theorem II
  (dimension exactly 2 under an irreducibility hypothesis on $\mathbb{Q}(\zeta_p)$), Theorem III
  (``$\tilde R^0_{\bar\rho}$ and $A_{\bar\rho}$ have dimension at least 2''). $A_{\bar\rho}$ is the shallow
  mod-$p$ algebra generated by $T_\ell$ and $S_\ell$, a limit over weights. Not in these theorems: ``$x,y$
  realisable as $T_q-a_q$'' and ``irreducible'' (as in [E] \texttt{cl:bk}).
  \textbf{Guards} (the source's standing hypotheses, L16-5 finding~1): $\operatorname{char}F=p$ and $\bar\rho$ modular
  on all three theorems, $p>3$ on Theorems~II and~III (``In \S\S 2--9 we assume $p>3$'').
\end{hypothesis}

\section{Bernoulli numbers and Iwasawa theory}

\begin{hypothesis}[Washington, Cor.~5.14 (Kummer congruences)]\label{lit:washington-5.14}
  \lean{GL2.Literature.Washington514} \leanok
  \textbf{Locator: verified, GTM 83 first edition, p.~61.} ``Suppose $m\equiv n\not\equiv0\pmod{p-1}$ are
  positive even integers. Then $B_m/m\equiv B_n/n\pmod p$. More generally, if $m$ and $n$ are positive even
  integers with $m\equiv n\pmod{(p-1)p^a}$ and $n\not\equiv0\pmod{p-1}$, then
  $(1-p^{m-1})B_m/m\equiv(1-p^{n-1})B_n/n\pmod{p^{a+1}}$.'' Sanity instance $p=5$, $(m,n)=(2,6)$ decided in Lean.
\end{hypothesis}

\begin{hypothesis}[Washington, Thm.~7.10 (derived)]\label{lit:washington-7.10}
  \lean{GL2.Literature.WashingtonIwasawaBernoulli} \leanok
  \textbf{Locator: verified, first edition, p.~123.} ``Let $\chi=\theta\psi$ be an even Dirichlet character
  ($\theta$ = first kind, $\psi$ = second kind) \dots\ Then $L_p(s,\chi)=f(\zeta_\psi(1+q_0)^s-1,\theta)$,''
  with $f\in\mathcal{O}_\theta[[T]]$ (Thm.~7.1, Prop.~7.6, p.~122). It is the interpolation formula, not a
  statement about $f(a-1+aY)$ as quoted in slot-15 WP-6. \textbf{The Lean structure is a derived instance}
  (Thms.~7.10, 7.1, 7.6 and 5.11 specialised to $\theta=\omega^{k_0}$, $\psi=1$, $s=1-k$), not the quoted theorem;
  hence its name \texttt{WashingtonIwasawaBernoulli} (L16-5 minor~8).
\end{hypothesis}

\section{Hecke operators on level-one forms (Serre, \emph{A course in arithmetic})}

\begin{hypothesis}[Serre, \emph{Course}, VII \S5.4, Thm.~7 and Cor.~1]\label{lit:serre-course-thm7}
  \lean{GL2.Literature.SerreCoursThm7} \leanok
  \textbf{Locator: verified, GTM 7, Ch.~VII, \S5.4, printed p.~102} (page image; the \emph{Cours} is in the library
  as \emph{A Course in Arithmetic}). ``Let $f(z)=\sum c(n)q^n$ be a modular form of weight $2k$, $k>0$, not identically
  zero. We assume that $f$ is an eigenfunction of all the $T(n)$, i.e.\ that there exists a complex number $\lambda(n)$
  such that $T(n)f=\lambda(n)f$ for all $n\ge1$. \textbf{Theorem 7.} a) The coefficient $c(1)$ of $q$ in $f$ is $\ne0$.
  b) If $f$ is normalized by the condition $c(1)=1$, then $c(n)=\lambda(n)$ for all $n>1$. \dots\ \textbf{Corollary 1.}
  Two modular forms of weight $2k$, $k>0$, which are eigenfunctions of the $T(n)$ with the same eigenvalues
  $\lambda(n)$, and which are normalized, coincide.''
  \textbf{Guards} (the source's standing hypotheses): $K=2k>0$ even, and formula (76) (Prop.~12) for the operators,
  which are interface data. Multiplicity one only: the \emph{existence} of an eigenbasis is not in the \emph{Course}
  (it is Diamond--Im \S3.6, \texttt{lit:diamond-im-eigen}, with \texttt{lit:serre-course-prop13}).
  Part b) is \emph{not} a field (L18-4; the field \texttt{coeff\_eq\_eigenvalue} was dropped): Serre's proof of b) uses only
  the one $n$ and \textbf{Prop.~12, Cor.~1} (``$\gamma(1)=c(n)$'', printed p.~102), i.e.\ formula (76) at $m=1$, the guard
  itself, so b) is proved from the guard (\texttt{inf:serre-thm7b}; L17-5 audit MINOR-1).
  Used by [C] \texttt{lem:heckerankone}, \texttt{prop:cohforms}(3).
\end{hypothesis}

\begin{lemma}[Serre, \emph{Course}, VII \S5.4, Thm.~7 b), from the guard]\label{inf:serre-thm7b}
  \lean{GL2.Literature.serre_coeff_eq_eigenvalue}
  \leanok
  \uses{lit:serre-course-thm7}
  Under formula (76), a normalised $f$ ($c(1)=1$) with $T(n)f=\lambda f$ for the one $n\ge1$ has $c(n)=\lambda$.
  \gtmodlit{formula (76) for the operators (Serre, Prop.~12, the guard of \texttt{lit:serre-course-thm7}), as the hypothesis
  \gtlithyp{hT} (\texttt{HeckeCoeffFormula})}
\end{lemma}
\begin{proof}
  \leanok
  (L18-4.) The coefficient of $q$ in $T(n)f$ is $\sum_{a\mid(n,1)}a^{K-1}c(n/a^2)=c(n)$ by (76) at $m=1$, and it is
  $\lambda c(1)=\lambda$.
\end{proof}

\begin{hypothesis}[Serre, \emph{Course}, VII \S5.4, Thm.~7 Cor.~2, (79)--(80)]\label{lit:serre-course-thm7-cor2}
  \lean{GL2.Literature.SerreCoursThm7Cor2} \leanok
  \textbf{Locator: verified, GTM 7, Ch.~VII, \S5.4, printed p.~102} (page image, L17-5 repairs). ``\textbf{Corollary 2.}
  Under the hypothesis of theorem 7, b): (79) $c(m)c(n)=c(mn)$ if $(m,n)=1$, (80) $c(p)c(p^n)=c(p^{n+1})+p^{2k-1}c(p^{n-1})$.
  Indeed the eigenvalues $\lambda(n)=c(n)$ satisfy the same identities (72) and (73) as the $T(n)$.''
  \textbf{Guards:} $K=2k>0$ even, formula (76); $m,n\ge1$, $p$ prime, $n\ge1$ in (80). The pinned variant of
  \texttt{lit:serre-course-prop11} on $M_K$. Used by [C] \texttt{thm:relations} (forms side), \texttt{lem:heckerankone}.
\end{hypothesis}

\begin{hypothesis}[Serre, \emph{Course}, VII \S5.5, Prop.~13]\label{lit:serre-course-prop13}
  \lean{GL2.Literature.SerreCoursProp13} \leanok
  \textbf{Locator: verified, GTM 7, Ch.~VII, \S5.5 a), printed p.~104} (page image, L17-5 repairs). ``Let $k$ be an integer
  $\ge2$. \textbf{Proposition 13.} The Eisenstein series $G_k$ is an eigenfunction of $T(n)$; the corresponding eigenvalue is
  $\sigma_{2k-1}(n)$ and the normalized eigenfunction is (86) $(-1)^k\frac{B_k}{4k}E_k=(-1)^k\frac{B_k}{4k}+\sum_{n=1}^\infty
  \sigma_{2k-1}(n)q^n$. The corresponding Dirichlet series is $\zeta(s)\zeta(s-2k+1)$.''
  \textbf{Guards:} $K=2k\ge4$ even, formula (76). Stated weaker: \emph{some} form with coefficients $\sigma_{K-1}(n)$, $n\ge1$,
  is an eigenform of every $T(n)$ with eigenvalue $\sigma_{K-1}(n)$ (no identification with $G_k$, no constant term). With
  $M_k=M^0_k\oplus\mathbf C.G_k$ (for $k\ge2$; \S3.2, printed p.~88, verified, not a structure) this is ``the Eisenstein
  system once'' of [C] \texttt{lem:heckerankone}.
\end{hypothesis}

\section{Eichler--Shimura and Hecke eigenspaces (Diamond--Im)}

\begin{hypothesis}[Diamond--Im, Thm.~12.2.2 and Prop.~12.4.10]\label{lit:diamond-im-es}
  \lean{GL2.Literature.DiamondImEichlerShimura} \leanok
  \textbf{Locator: verified, Diamond--Im, \emph{Modular forms and modular curves}, CMS Conf.\ Proc.\ 17 (1995), printed
  pp.~102, 107--108, 115, 117} (camera scan, 150 dpi, quarter-cropped; L17-5 audit MAJOR-1, read in the repairs).
  \S12.1: ``Let $k$ be an integer and $\Gamma$ a congruence subgroup of $\mathrm{SL}_2(\mathbf Z)$.'' \S12.2: ``We let
  $V_k=\mathrm{Symm}^{k-2}_{\mathbf C}(\mathbf C^2)$ \dots\ we obtain a $\mathbf C$-linear map $\beta\colon\mathcal M_k(\Gamma)\oplus
  \overline{\mathcal S}_k(\Gamma)\to H^1(\Gamma,V_k)$ which restricts to $\beta_p\colon\mathcal S_k(\Gamma)\oplus\overline{\mathcal S}_k(\Gamma)
  \to H^1_p(\Gamma,V_k)$. \textbf{Theorem 12.2.2.} $\beta$ and $\beta_p$ are isomorphisms.'' \S12.4: ``$\Gamma=\Gamma_1(N)$ or
  $\Gamma_0(N)$. Let $\Delta$ denote the corresponding semigroup $\Delta_1(N)$ or $\Delta_0(N)$ \dots\ \textbf{Proposition 12.4.10.}
  The Eichler--Shimura isomorphisms $\beta$ and $\beta_p$ of Theorem 12.2.2 are compatible with the action of $R(\Gamma,\Delta)$.''
  \textbf{Guards:} $k\ge2$ (\S12.2's weight), $N\ge1$. Stated weaker: \emph{some} Hecke-equivariant isomorphisms exist.
  Interface data for $\Gamma=\Gamma_0(N)$ or $\Gamma_1(N)$; at level one $\Gamma=\mathrm{SL}_2(\mathbf Z)$ and $V_k=\mathrm{Sym}^{n}\otimes\mathbf C$,
  $k=n+2$. The p.~117 hypothesis ``$\Gamma=\Gamma_1(N)$ with $N>4$'' governs only the sheaf description (12.4.6).
  Used by \texttt{lit:eichler-shimura-levelone}, \texttt{lit:forms-same-systems}.
\end{hypothesis}

\begin{hypothesis}[Diamond--Im, \S3.5--\S3.6: eigenspaces]\label{lit:diamond-im-eigen}
  \lean{GL2.Literature.eigenC, GL2.Literature.DiamondImEigenspaces} \leanok
  \textbf{Locator: verified, Diamond--Im, printed pp.~54--55} (L17-5 repairs). ``Let us next observe that every
  $\mathbf T_N$-eigenspace is (at most) one-dimensional. \dots\ Thus, if two forms of type $(k,N,\varepsilon)$ are common
  eigenforms of $T_n$ for all $n$ with the same system $\{\lambda_n\}$ of eigenvalues then one is a scalar multiple of the
  other.'' and ``\dots\ we deduce that there is an orthogonal decomposition $\mathcal S_k(\Gamma_1(N))=\bigoplus_\varepsilon
  \mathcal S_k(N,\varepsilon)$ \dots, and that each $\mathcal S_k(N,\varepsilon)$ decomposes orthogonally into a direct sum of
  $\mathbf T^{(N)}$-eigenspaces [Miy2, Theorem 4.5.4].''
  \textbf{Guard:} $N\ge1$. Orthogonality not recorded. At $N=1$, $\mathbf T^{(1)}$ is the algebra of all $T_n$: the
  eigenbasis of $S_k(\mathrm{SL}_2(\mathbf Z))$. Used by \texttt{lit:eichler-shimura-levelone}, \texttt{lit:forms-same-systems}.
\end{hypothesis}

\begin{hypothesis}[Serre, \emph{Course}, VII \S5.3, Prop.~11, (72)--(73)]\label{lit:serre-course-prop11}
  \lean{GL2.Literature.SerreCoursProp11} \leanok
  \textbf{Locator: verified, GTM 7, Ch.~VII, \S5.3, printed p.~101.} ``\textbf{Proposition 11.} The function $T(n)f$ is
  weakly modular of weight $2k$ \dots\ We have: (72) $T(m)T(n)f=T(mn)$ if $(m,n)=1$, (73)
  $T(p)T(p^n)f=T(p^{n+1})f+p^{2k-1}T(p^{n-1})f$, if $p$ is prime, $n\ge1$.''
  \textbf{Guards:} $K=2k>0$ even; they are vacuous on the free interface data $T:\mathbb N\to\operatorname{End}M$
  (L17-5 audit MINOR-2): nothing ties $T$ to Serre's $T(n)$, the L2 layer must instantiate it; the pinned variant on
  $M_K$ is \texttt{lit:serre-course-thm7-cor2}. Numerical sanity
  ($\sigma_3(2)^2=\sigma_3(4)+2^3\sigma_3(1)$) decided in Lean. Used by [C] \texttt{thm:relations}.
\end{hypothesis}

\section{Hida theory, Emerton, Wiles}

\begin{hypothesis}[Hida (ens), Thm.~3.1]\label{lit:hida-ens-3.1}
  \lean{GL2.Literature.HidaEns31} \leanok
  \textbf{Locator: verified, Ann.\ Sci.\ \'ENS 19 (1986), p.~250.} ``Assume that $p\ge5$ and that $N$ is divisible
  by $p$. Then the ordinary parts $\mathscr{H}_0(N;\mathcal{O}_K)$ and $\mathfrak{h}_0(N;\mathcal{O}_K)$ \dots\ are
  free of finite rank over $\Lambda_K$. Moreover \dots\ $\operatorname{rank}_{\Lambda_K}(\mathfrak{h}_0(N,a;\mathcal{O}_K))
  =\operatorname{rank}_{\mathbb{Z}_p}(\mathscr{S}^0_k(\Phi;\mathbb{Z}_p))$ if $k\ge3$ and $k\equiv a\bmod p-1$.''
  \textbf{Guards} (L16-5 finding~1): $p\ge5$ and $p\mid N$, the level $N$ a parameter of the structure.
\end{hypothesis}

\begin{hypothesis}[Hida (ens), Cor.~3.2]\label{lit:hida-ens-3.2}
  \lean{GL2.Literature.HidaEns32} \leanok
  \textbf{Locator: verified, p.~250.} ``(3.3 b) $\mathfrak{h}_0(N,a;\mathcal{O}_K)/P_k\mathfrak{h}_0(N,a;\mathcal{O}_K)
  \cong\mathfrak{h}^0_k(\Gamma_1(N_0p),a;\mathcal{O}_K)$ if $k\ge2$.''
  \textbf{Guards} (L16-5 finding~1): $p\ge5$ and $p\mid N$ (Thm.~3.1's assumption). The Lean adds $k\equiv a\bmod p-1$,
  which (3.3\,b) does not have: a deliberate restriction to the class its consumer uses (a weakening; L16-5 minor~3).
\end{hypothesis}

\begin{hypothesis}[Hida (invent), Thm.~3.1(ii),(iii)]\label{lit:hida-inv-3.1}
  \lean{GL2.Literature.HidaInvent31} \leanok
  \textbf{Locator: verified, Invent.\ Math.\ 85 (1986), p.~562.} ``(ii) The $\Lambda$-modules $V^0$ and $W^0$ are
  free of finite rank over $\Lambda$. (iii) $\operatorname{rank}_\Lambda(V^0)=2\cdot\operatorname{rank}_\Lambda(\mathfrak{h}^0(N;\mathbb{Z}_p))$
  and $\operatorname{rank}_\Lambda(W^0)=\operatorname{rank}_\Lambda(V^0)+\tfrac12\varphi(p)\sum_{0<t|N}\varphi(t)\varphi(N/t)$.''
  \textbf{Guards} (L16-5 finding~1): $p\ge5$ and $(N,p)=1$ (\S0, p.~545: ``Assume throughout the paper that $p\ge5$'',
  ``$N$ prime to $p$'').
\end{hypothesis}

\begin{hypothesis}[Hida (invent), Prop.~4.7]\label{lit:hida-inv-4.7}
  \lean{GL2.Literature.HidaInvent47} \leanok
  \textbf{Locator: verified, p.~572.} ``For each integer $r>0$ and for each $n>0$, the restriction map combined
  with $e$ gives isomorphisms: $e_0H^1(\Gamma_1(N),L_n(\mathbb{Z}_p))\simeq eH^1(\Phi_r,L_n(\mathbb{Z}_p))$,
  $e_0H^1_P(\Gamma_1(N),L_n(\mathbb{Z}_p))\simeq eH^1_P(\Phi_r,L_n(\mathbb{Z}_p))$.'' Thm.~6.6 (p.~590) is
  the reduction-mod-$p^m$ companion (\texttt{lit:hida-inv-6.6}). The Lean structure states one isomorphism over one interface
  datum and is instantiated once per reading ($H^1$, resp.\ $H^1_P$; L16-5 minor~4).
\end{hypothesis}

\begin{hypothesis}[Hida (invent), Thm.~6.6 (reduction mod $p^m$)]\label{lit:hida-inv-6.6}
  \lean{GL2.Literature.HidaInvent66} \leanok
  \textbf{Locator: verified, Invent.\ Math.\ 85 (1986), p.~590} (page image). ``\textbf{Theorem 6.6.} For each $m>0$, there
  is a canonical isomorphism: $eH^1_P(\Phi_1,L_n(\mathbb Z_p))\otimes_{\mathbb Z_p}\mathbb Z/p^m\mathbb Z\simeq
  eH^1_P(\Phi_1,L_n(\mathbb Z/p^m\mathbb Z))$. Moreover, if $n\equiv a\bmod p-1$, we have an isomorphism of
  $\mathfrak h^0_2(Np^\infty;\mathbb Z_p)$-modules: \dots'' Only the first sentence is a field. \textbf{Guards:} $p\ge5$,
  $(N,p)=1$ (\S0), $m>0$, $n>0$. ``Canonical'' (naturality) is not recorded; the Hecke-module sentence is not stated.
\end{hypothesis}

\begin{hypothesis}[Emerton, Eisenstein ideal in Hida's algebra]\label{lit:emerton}
  \lean{GL2.Literature.EmertonThm} \leanok
  \textbf{Locator: verified, p.~2 of the library preprint.} ``Suppose that $i\not\equiv0\pmod{p-1}$. Then the pullback
  of $I^0_{\mathrm{Eis}}$ via the natural map $\Lambda_{(i)}\to\mathbf{T}^{\mathrm{ord}}_{\mathrm{par},(i)}$ is equal to the
  principal ideal of $\Lambda_{(i)}$ generated by the Kubota--Leopoldt $p$-adic $L$-function $L_p(\omega^i,1-s)$,
  and the consequent injection \dots\ is an isomorphism.''
  \textbf{Guards} (L16-5 finding~1): $p\ge5$ (the paper's standing hypothesis) and $i\not\equiv0\pmod{p-1}$, the component
  index $i$ a parameter of the structure.
\end{hypothesis}

\begin{hypothesis}[Wiles, Thm.~2.1.4]\label{lit:wiles-2.1.4}
  \lean{GL2.Literature.Wiles214} \leanok
  \textbf{Locator: verified, Invent.\ Math.\ 94 (1988), p.~561.} ``If $\mathbf f$ is ordinary at $\lambda$ and
  $\mathfrak p\mid N\lambda$ then the restriction of $\rho_\lambda$ to $D_{\mathfrak p}$ has the form
  $\rho_\lambda|_{D_{\mathfrak p}}\approx\left(\begin{smallmatrix}\varepsilon_1&*\\0&\varepsilon_2\end{smallmatrix}\right)$
  where $\varepsilon_2$ is unramified and $\varepsilon_2(\operatorname{Frob}\mathfrak p)=\alpha(\mathfrak p,\mathbf f)$ is a unit
  root of $x^2-c(\mathfrak p,\mathbf f)x+\psi(\mathfrak p)N\mathfrak p^{k-1}=0$.''
\end{hypothesis}

\section{The Eisenstein ideal}

\begin{hypothesis}[Mazur, II (9.7)]\label{lit:mazur-9.7}
  \lean{GL2.Literature.Mazur97} \leanok
  \textbf{Locator: verified, Publ.\ IH\'ES 47 (1977), p.~96.} ``$\mathbb T/\mathfrak J=\mathbb Z/n\mathbb Z$; the Eisenstein
  ideal $\mathfrak J$ contains the integer $n$,'' $n=\operatorname{numerator}((N-1)/12)$. FLT attests the weight-2
  statement (survey, F-A). Sanity instance $N=5$ proved.
\end{hypothesis}

\paragraph{Not in the register (reason).} Shimura (book, Ch.~3 and Ch.~8): held since 2026-10-05 and read in R20-L
(\texttt{docs/implementation/slot21/R20-L\_report.md}); the Hecke-ring locators below are verified there. Eichler--Shimura is F-A
upstream, and since the L17-5 repairs it is no longer needed for \texttt{lit:eichler-shimura-levelone}
and \texttt{lit:forms-same-systems} (Part G), whose sources are Diamond--Im \S3.5--3.6, \S12 and Serre VII \S3.2, \S5 (above).
Sturm LNM 1240 (received 2026-10-05): Thm.~1, p.~276 verified, see \texttt{lit:sturm}. Jochnowitz, TAMS 270 (1982) 253--267
(received 2026-10-05): Lemma~4.1(b), p.~257 and \S7, pp.~261--264, read (the ``Thms.~6.3/6.6, p.~280'' and ``Thm.~4.1, p.~276''
citations are to the other Jochnowitz paper, TAMS 270, 269--285). Merel 1996 (LN5), Wake--Wang-Erickson (LN6), Katz 1977,
Serre LNM 350: not in the library; locators stay VERIFY.
Slot 20 (L20-2, L20-5): Shimura Thm.~3.24, Prop.~3.8, 3.16, 3.17, Brown III.9 and Ash--Stevens \S1 enter the register below as
\emph{transported} statements on the model; their locators were read in R20-L (2026-10-05), see each row.

\section{The Hecke-ring relations on the model (Shimura; transported) --- slot 20, L20-2, branch H}

% L20-0 §3.2: three rows, each stated on the model's operators (GL2.M4.Tdc / Tprim / TdcOf) in the project's notation, NOT in
% Shimura's T(a,d) / T(m) (R20-L has since read the scan and confirmed Shimura's notation, Thm. 3.24(1), (2)).
\providecommand{\VERIFYloc}[1]{\textbf{Locator: #1.}}

\paragraph{Transport.} Shimura's identities live in the abstract Hecke ring $R(\Gamma,\Delta)$; on $H^n$ they need the
multiplicativity of the action of $R(\Gamma,\Delta)$ on cohomology (Ash--Stevens, Crelle 365 (1986), \S1; Rhie--Whaples).
Neither the ring nor the action is in mathlib \texttt{db584cd}, so each row states the \emph{transported} identity directly
on the operators of \ref{def:m4-hecke}; there is no separate \texttt{lit:hecke-action-mult} row (\ref{rmk:m4-g5-headers}).
\textbf{Guards} (all three): $6\in R^\times$; $\rho(pI)=c$ for the scalar $z=pI$; $M,N>0$ coprime in the second row.
\textbf{Numerical check against the C++} (L20-0 \S3.2(b); quoted, not re-run): \texttt{HeckeRelationsValidation}, oracle O9
([C] \texttt{obl:algebra}, all $\ell\ge5$; \texttt{Hecke/README.md}): O9.2 ``$T_4=T_2^2-2^{n+1}I$, FULL operator, independent
route'' with O9.5 ``decomposition witness $T_4^{\mathrm{full}}=T'_4+2^nI$'' are together \texttt{lit:hecke-Tp-prim} at $p=2$,
$\mu=0$ ($T_2^2=T'_4+(2+1)2^n$); O9.1 (commutativity of $T_2$, $T_3$) and O9.3 ($T_6=T_2T_3$) are instances of
\texttt{lit:hecke-commute} and \texttt{lit:hecke-coprime}.

\paragraph{The nebentypus identities (no separate row; for R20-L).} \ref{thm:nebrelations-model} takes the double-coset
products of the primitive operators \emph{with character} in $\mathcal H(\Gamma_0(N),\Delta_0(N))$, transported to $H^n$, as the
inline hypotheses \texttt{hTpchi0} ($\mu=0$: $T_p^2=T_{\Delta^{\mathrm{prim}}_{p^2}}+(p+1)\chi(p)c$, $p\nmid N$), \texttt{hTpchi}
($\mu\ge1$: $T_pT_{\Delta^{\mathrm{prim}}_{p^{\mu+1}}}=T_{\Delta^{\mathrm{prim}}_{p^{\mu+2}}}+p\,\chi(p)c\,T_{\Delta^{\mathrm{prim}}_{p^\mu}}$)
and \texttt{hUchi} ($p\mid N$: $U_pU_{p^\mu}=U_{p^{\mu+1}}$), the \texttt{hrank} precedent for inline binders.
\VERIFYloc{Shimura \S3.3: $\Gamma'$ (3.3.2), p.~67, is $\Gamma_0(N)$ and $\Delta'$ (3.3.3), p.~68, is $\Delta_0(N)$ at $t=1$,
$\mathfrak h=(\mathbb Z/N)^\times$; \texttt{hTpchi0} / \texttt{hTpchi} are Thm.~3.24(5) with Prop.~3.17 carried to
$R(\Gamma_0(N),\Delta_0(N))$ by Prop.~3.31, p.~71, Shimura's ring carrying no character, which enters only through
$T(p,p)\mapsto\chi(p)c$; \texttt{hUchi} is Thm.~3.34(3) with Prop.~3.33; the twist is Ash--Stevens Lemma~1.1.5, p.~196 ---
verified, R20-L}
(F3-20 MINOR 9; read in R20-L, 2026-10-05.)

\begin{hypothesis}[Shimura, Thm.~3.24(5) / Prop.~3.17: $T_{\Delta^{\mathrm{prim}}_p}T_{\Delta^{\mathrm{prim}}_{p^{\mu+1}}}$, transported]\label{lit:hecke-Tp-prim}
  \lean{GL2.Literature.HeckeTpPrim} \leanok
  \VERIFYloc{Shimura, \emph{Introduction to the arithmetic theory of automorphic functions}, Thm.~3.24(5), p.~63, with
  Prop.~3.17, p.~58 ($T(p,p^k)=T(p,p)T(1,p^{k-1})$); transported by Ash--Stevens \S1.1, p.~194 --- verified, R20-L}
  No quotation. Stated (\texttt{hTp}): for every family of $\tau$s, on $H^n$,
  $\mathtt{Tprim}_1^2=\mathtt{Tprim}_2+(p+1)c$ and, for $\mu\ge1$,
  $\mathtt{Tprim}_1\,\mathtt{Tprim}_{\mu+1}=\mathtt{Tprim}_{\mu+2}+p\,c\,\mathtt{Tprim}_\mu$ (the product of the primitive double
  cosets of determinant $p$ and $p^{\mu+1}$, with $T(p,p)$ acting by $c$).
\end{hypothesis}

\begin{hypothesis}[Shimura, Prop.~3.16, the coprime clause, transported]\label{lit:hecke-coprime}
  \lean{GL2.Literature.HeckeCoprime} \leanok
  \VERIFYloc{Shimura, Prop.~3.16, p.~57 (Thm.~3.24(3) is the full-operator form); transported by Ash--Stevens \S1.1, p.~194
  --- verified, R20-L}
  No quotation. Stated (\texttt{hcop}): for coprime $M,N>0$,
  $T_{\Gamma\mathrm{diag}(1,M)\Gamma}T_{\Gamma\mathrm{diag}(1,N)\Gamma}=T_{\Gamma\mathrm{diag}(1,MN)\Gamma}$ on $H^n$, for the choice-free
  operators \texttt{TdcOf} (one coset datum and one $\tau$ each; every other datum by \texttt{TdcOf\_eq}).
\end{hypothesis}

\begin{hypothesis}[Shimura, Prop.~3.8: commutativity, transported]\label{lit:hecke-commute}
  \lean{GL2.Literature.HeckeCommute} \leanok
  \VERIFYloc{Shimura, Prop.~3.8, p.~54, as applied in \S3.2, p.~56 ($\det>0$); $\det<0$ by the diagonal-representative remark
  below; transported by Ash--Stevens \S1.1, p.~194 (whose Hecke algebras are assumed commutative, p.~193) --- verified, R20-L}
  No quotation. Stated (\texttt{hcomm}): $T_{\Gamma g\Gamma}$ and $T_{\Gamma g'\Gamma}$ commute on $H^n$ for all $g,g'\in\Delta$
  (choice-free \texttt{TdcOf}). L20-0 \S3.2(c): this is \emph{not} implied by the other two rows for an arbitrary $g$. F3-20 (MINOR
  4): it is [C] \texttt{obl:algebra}'s ``whence the $\{T_N\}$ commute'' (the image of a commutative $\mathcal H(\Delta,\Gamma)$), of
  which \texttt{thm:relations}(1)'s first clause is an instance; it is stronger than that clause as literally written. $\Delta$
  here includes $\det<0$; Prop.~3.8 applies since every double coset has a symmetric (diagonal) representative (Smith form, the
  sign of $\det$ moved onto $d$ by $\mathrm{diag}(1,-1)$).
\end{hypothesis}

\begin{theorem}[Proved (M6-A): Brown III.9.5, Ash--Stevens \S1.1: $T_{\Gamma g\Gamma}=\mathrm{cor}\circ(c_g)_*\circ\mathrm{res}$, on the model]\label{lit:hecke-classical}
  % slot 20, L20-5 (M4-f, branch H; L20-0 §3.3).  Lean: lean/GL2/Literature/HeckeClassical.lean.
  % slot 23 (M6-A): no longer a hypothesis -- proved (GL2.M6.heckeClassical_holds, thm:m6-hcl); the label is kept.
  \lean{GL2.Literature.HeckeClassical, GL2.M6.heckeClassical_holds} \leanok
  \VERIFYloc{Brown, \emph{Cohomology of groups} (GTM 87), Prop.~III.9.5(ii), p.~82 ($\mathrm{cor}\circ\mathrm{res}=(G:H)$) and
  (iii), pp.~82--83 (the double-coset formula; not \S10); $c(g)^*$ is Brown \S8, p.~80; Weibel 6.7.17, p.~194, is a second leg for
  $\mathrm{cor}\circ\mathrm{res}=[G:H]$ --- verified. \textbf{Ash--Stevens, Crelle 365 (1986), \S1.1, pp.~193--196, is NOT AS CITED}:
  it gives $T_g$ by the explicit cocycle formula (p.~194) and states res / cor compatibility (Lemma~1.1.3), but never writes
  $T_g=\mathrm{cor}\circ c_g\circ\mathrm{res}$. The identity is standard and derivable from Ash--Stevens p.~194 with Brown \S9(D);
  Kuga--Parry--Sah (Ash--Stevens' [20]) was not read --- R20-L}
  No quotation. Stated (\texttt{hcl}): for $6\in R^\times$, $g\in\Delta$ and $n$, there is an $R$-linear
  $\mathrm{cor}\colon H^n(\Gamma'';V_L)\to H^n(\mathrm{SL}_2(\mathbb Z);V_L)$ with $\mathrm{cor}\circ\mathrm{res}=[\Gamma:\Gamma'']$ and
  $\Theta\circ\texttt{TdcOf}~g=\mathrm{cor}\circ(c_g)_*\circ\mathrm{res}\circ\Theta$ (the objects of \ref{def:m4-classical}).
  $\mathrm{cor}$ is hypothesised because mathlib \texttt{db584cd} has no corestriction. \textbf{Weak form} (F3-20 MAJOR 1):
  $\mathrm{cor}\circ\mathrm{res}=[\Gamma:\Gamma'']$ does not determine $\mathrm{cor}$, so the row does not identify it with the
  corestriction; [C] \texttt{thm:classical} stays [C]-only. \textbf{Guards:} $6\in R^\times$, $g\in\Delta$.
  The locators of this row are Brown's and Ash--Stevens', not Shimura's; R20-L read them (2026-10-05): Brown's are verified,
  the Ash--Stevens leg is not as cited (above), so \texttt{hcl} rests on the derivation, not on a quotation. Weibel 6.7.17 (owner
  transcription): ``The composite $\mathrm{cor}^G_H\circ\mathrm{tr}$ is multiplication by the index $[G:H]$ on $H_*(G;A)$. Similarly,
  the composite $\mathrm{tr}\circ\mathrm{res}$ is multiplication by $[G:H]$ on $H^*(G;A)$''; Weibel's cohomological $\mathrm{tr}$ is
  our $\mathrm{cor}$, so our $\mathrm{cor}\circ\mathrm{res}=[\Gamma:\Gamma'']$ is the second sentence.
  \emph{Slot 23 (M6-A): proved}, with $\mathrm{cor}$ \emph{the} corestriction (\ref{thm:m6-hcl}); this row is no longer a literature
  hypothesis.
\end{theorem}
\begin{proof}
  \leanok
  \ref{thm:m6-hcl}.
\end{proof}
